% Chapter 3.
\chapter{A hybrid estimator}
\label{ch:estimation}

\section{Learning the residual}
\label{sec:residual}

The ageing law of \cref{eq:ageing} captures the broad trend of every site
but leaves a residual of up to two per cent in state of health. Rather
than adding terms to the physical model until it fits, the hybrid
estimator treats the law as a prior mean and learns what it misses from
data \parencite{rasmussen2006gp}. Let \(m(\mathbf{x}_t)\) denote the prediction of the ageing law for
the operating history \(\mathbf{x}_t\) up to time \(t\). The estimator
models the true state of health as
\begin{equation}
  \mathrm{SoH}(t) = m(\mathbf{x}_t) + f(\mathbf{z}_t) + \varepsilon,
  \qquad f \sim \mathcal{GP}\bigl(0, k(\mathbf{z}, \mathbf{z}')\bigr),
  \label{eq:hybrid}
\end{equation}
where \(\mathbf{z}_t\) is a vector of features extracted from recent
partial charge segments and \(\varepsilon\) is the measurement noise of
the reference test.

\begin{definition}{Partial charge feature}{feature}
  A \emph{partial charge feature} is a statistic of the voltage curve
  recorded while the battery charges continuously across a fixed
  window of state of charge, here 40 to 60\,\%. The features used are the
  area under the incremental capacity curve \(\mathrm{d}Q/\mathrm{d}V\)
  within the window, the voltage at which its peak occurs, and the
  equivalent series resistance at the start of the segment.
\end{definition}

Frequency-response duty rarely charges a battery through a full window,
but over three years each site accumulates several hundred qualifying
segments, enough to estimate the features weekly. The kernel \(k\) is a
sum of a squared-exponential term over the features and a linear term
over time, which lets the correction drift slowly while responding to
changes in electrode behaviour.

\section{Evaluation on held-out sites}
\label{sec:evaluation}

The estimator was evaluated by leaving one site out at a time: the
ageing law and the Gaussian process were fitted on ten sites and used to
track the eleventh from commissioning onward, with access to its
telemetry but none of its reference tests. \Cref{tab:errors} compares
the hybrid estimator with its two components and with the estimate
reported by each site's battery management system.

\begin{table}[tb]
  \caption[Estimation error on held-out sites]{Estimation error on
    held-out sites, in percentage points of state of health. Coverage is
    the share of reference tests inside the 95\,\% predictive interval.}
  \label{tab:errors}
  \centering
  \begin{tabular}{@{}l r r r@{}}
    \toprule
    Estimator & RMSE & Max.\ error & Coverage \\
    \midrule
    Battery management system  & 3.41 & 7.20 & --- \\
    Ageing law alone           & 1.92 & 4.05 & --- \\
    Gaussian process alone     & 2.37 & 6.81 & 71\,\% \\
    \textbf{Hybrid estimator}  & \textbf{0.88} & \textbf{2.14} & \textbf{94\,\%} \\
    \bottomrule
  \end{tabular}
\end{table}

The hybrid estimator halves the error of the ageing law and divides that
of the management system by almost four. More importantly, its
uncertainty is honest: 94\,\% of the held-out reference tests fall inside
the nominal 95\,\% interval. The Gaussian process alone, lacking a
physical mean, extrapolates poorly beyond the ages seen in training and
is overconfident there.

\begin{figure}[tb]
  \marginfigure{%
  \begin{tikzpicture}
    \begin{axis}[
        width=0.98\linewidth, height=0.6\linewidth,
        xmin=0, xmax=36, ymin=0.90, ymax=1.005,
        xtick={0,6,...,36}, ytick={0.90,0.925,0.95,0.975,1.0},
        yticklabel style={/pgf/number format/.cd, fixed, precision=3},
        xlabel={Months in service}, ylabel={State of health},
        legend pos=south west,
      ]
      \addplot[name path=upper, draw=none, domain=0:36, samples=60]
        {1 - 0.0150*sqrt(x) - 0.00090*x + 0.002 + 0.00035*x};
      \addplot[name path=lower, draw=none, domain=0:36, samples=60]
        {1 - 0.0150*sqrt(x) - 0.00090*x - 0.002 - 0.00035*x};
      \addplot[accent!15] fill between[of=upper and lower];
      \addlegendentry{95\,\% interval}
      \addplot[accent, line width=1pt, domain=0:36, samples=60]
        {1 - 0.0150*sqrt(x) - 0.00090*x};
      \addlegendentry{Hybrid estimate}
      \addplot[black!55, line width=0.8pt, densely dashed, domain=0:36, samples=60]
        {1 - 0.0085*sqrt(x) - 0.00062*x};
      \addlegendentry{Management system}
      \addplot[only marks, mark=*, mark size=1.5pt, black]
        coordinates {(6,0.9578) (12,0.9367) (19,0.9177) (25,0.9017) (33,0.8826)};
      \addlegendentry{Reference test}
    \end{axis}
  \end{tikzpicture}}{%
  \caption[Tracking a held-out site]{Tracking site NL-04 without any of
    its reference tests. The management system drifts high; the hybrid
    estimate stays within its interval.}
  \label{fig:tracking}}
\end{figure}

\begin{keyresult}{Accuracy without site-specific tests}{accuracy}
  Trained on other sites only, the hybrid estimator tracks the state of
  health of a new site to within 0.88 percentage points RMSE over three
  years, with a calibrated 95\,\% interval. No reference test of the
  target site is needed after commissioning.
\end{keyresult}

\section{Field trial}
\label{sec:trial}

Operators keep a battery between conservative state-of-charge limits,
typically 10 and 90\,\%, partly because they do not trust their estimate
of how much capacity remains. If the estimate is known to within a
percentage point, the lower limit can follow the true capacity rather
than a worst case. From February to October 2026, site NL-01 ran with
its lower limit set by the hybrid estimate minus twice its standard
deviation, while a twin container on the same site kept the fixed limit.

Over nine months the trial container delivered 6.4\,\% more energy into
the frequency reserve market. Its reference test at the end of the trial
showed a state of health of 0.947, against 0.949 for the control, a
difference well within the test's repeatability of \(\pm 0.004\). The
estimate itself, made before the test, was 0.944.
